\documentclass[10pt,a4paper]{amsart}
\usepackage{amsmath,amssymb,amsthm,graphicx}
% Retained minimal preamble for the Stacks Project chapter more-morphisms.tex (More on Morphisms).
% Verbatim excerpt of the chapter preamble: Xy-pic setup (native preamble lines 14-18) and the native
% theorem declarations with the mathematical macros (native preamble lines 154-217). Copied mechanically;
% nothing added, removed or redefined. The wrapper class and the amsmath/amssymb/amsthm/graphicx base
% packages are supplied by the assembler, so the class/externaldocument part of the native preamble is omitted.
% For commutative diagrams we use Xy-pic
\usepackage[all]{xy}

% We use 2cell for 2-commutative diagrams.
\xyoption{2cell}
% Theorem environments.
%
\theoremstyle{plain}
\newtheorem{theorem}[subsection]{Theorem}
\newtheorem{proposition}[subsection]{Proposition}
\newtheorem{lemma}[subsection]{Lemma}

\theoremstyle{definition}
\newtheorem{definition}[subsection]{Definition}
\newtheorem{example}[subsection]{Example}
\newtheorem{exercise}[subsection]{Exercise}
\newtheorem{situation}[subsection]{Situation}

\theoremstyle{remark}
\newtheorem{remark}[subsection]{Remark}
\newtheorem{remarks}[subsection]{Remarks}

\numberwithin{equation}{subsection}

% Macros
%
\def\lim{\mathop{\mathrm{lim}}\nolimits}
\def\colim{\mathop{\mathrm{colim}}\nolimits}
\def\Spec{\mathop{\mathrm{Spec}}}
\def\Hom{\mathop{\mathrm{Hom}}\nolimits}
\def\Ext{\mathop{\mathrm{Ext}}\nolimits}
\def\SheafHom{\mathop{\mathcal{H}\!\mathit{om}}\nolimits}
\def\SheafExt{\mathop{\mathcal{E}\!\mathit{xt}}\nolimits}
\def\Sch{\mathit{Sch}}
\def\Mor{\mathop{\mathrm{Mor}}\nolimits}
\def\Ob{\mathop{\mathrm{Ob}}\nolimits}
\def\Sh{\mathop{\mathit{Sh}}\nolimits}
\def\NL{\mathop{N\!L}\nolimits}
\def\CH{\mathop{\mathrm{CH}}\nolimits}
\def\proetale{{pro\text{-}\acute{e}tale}}
\def\etale{{\acute{e}tale}}
\def\QCoh{\mathit{QCoh}}
\def\Ker{\mathop{\mathrm{Ker}}}
\def\Im{\mathop{\mathrm{Im}}}
\def\Coker{\mathop{\mathrm{Coker}}}
\def\Coim{\mathop{\mathrm{Coim}}}

% Boxtimes
%
\DeclareMathSymbol{\boxtimes}{\mathbin}{AMSa}{"02}

%
% Macros for moduli stacks/spaces
%
\def\QCohstack{\mathcal{QC}\!\mathit{oh}}
\def\Cohstack{\mathcal{C}\!\mathit{oh}}
\def\Spacesstack{\mathcal{S}\!\mathit{paces}}
\def\Quotfunctor{\mathrm{Quot}}
\def\Hilbfunctor{\mathrm{Hilb}}
\def\Curvesstack{\mathcal{C}\!\mathit{urves}}
\def\Polarizedstack{\mathcal{P}\!\mathit{olarized}}
\def\Complexesstack{\mathcal{C}\!\mathit{omplexes}}
% \Pic is the operator that assigns to X its picard group, usage \Pic(X)
% \Picardstack_{X/B} denotes the Picard stack of X over B
% \Picardfunctor_{X/B} denotes the Picard functor of X over B
\def\Pic{\mathop{\mathrm{Pic}}\nolimits}
\def\Picardstack{\mathcal{P}\!\mathit{ic}}
\def\Picardfunctor{\mathrm{Pic}}
\def\Deformationcategory{\mathcal{D}\!\mathit{ef}}
\usepackage{cleveref}
\crefname{originalreference}{original reference}{original references}
\Crefname{originalreference}{Original reference}{Original references}
\makeatletter
\newlabel{algebra-lemma-qc-open}{{00F6}{source}{}{sourcepointer-1}{}}
\newlabel{algebra-lemma-qc-open@cref}{{[originalreference][0][]00F6}{[1][0][]source}{}{}{}}
\makeatother
\crefname{originallemma}{original lemma}{original lemmas}
\Crefname{originallemma}{Original lemma}{Original lemmas}
\makeatletter
\newlabel{lemma-stratify-flat-fp-lqf-universally-bounded}{{07RY}{source}{}{sourcepointer-2}{}}
\newlabel{lemma-stratify-flat-fp-lqf-universally-bounded@cref}{{[originallemma][0][]07RY}{[1][0][]source}{}{}{}}
\makeatother
\crefname{originallemma}{original lemma}{original lemmas}
\Crefname{originallemma}{Original lemma}{Original lemmas}
\makeatletter
\newlabel{lemma-stratify-flat-fp-qf}{{07RZ}{source}{}{sourcepointer-3}{}}
\newlabel{lemma-stratify-flat-fp-qf@cref}{{[originallemma][0][]07RZ}{[1][0][]source}{}{}{}}
\makeatother
\crefname{originalreference}{original reference}{original references}
\Crefname{originalreference}{Original reference}{Original references}
\makeatletter
\newlabel{modules-definition-support}{{01AT}{source}{}{sourcepointer-4}{}}
\newlabel{modules-definition-support@cref}{{[originalreference][0][]01AT}{[1][0][]source}{}{}{}}
\makeatother
\title{Annihilators of sections go down along flat locally quasi-finite morphisms}
\author{The Stacks Project Authors}
\date{stacks-project:more-morphisms.tex:local-snapshot-20261011; GNU FDL 1.2 or later}
\begin{document}
\maketitle
\noindent\textit{Editorial note.} Counted span is the author's own proof environment, source lines 12932-12976, immediately following the statement at 12913-12930.\par\medskip
\noindent\textit{Editorial note.} The proof constructs psi and the map A[I]\textasciicircum{}\{r\} -> B[I] explicitly; both are displayed inside the span.\par\medskip
\noindent\textit{Editorial note.} Three cross-chapter/same-chapter pointers (Algebra, Modules) are recorded in assembly.external\_references.\par\medskip
\noindent\textit{Allowed standard prerequisites (editorial).}\begin{itemize}
\item Stratification by fibre degree in the quasi-finite case. chapter lemma lemma-stratify-flat-fp-qf, cited in the proof. Provides the stratification by fibre degree used to reduce to a finite locally free piece.
\item Uniformly bounded stratification. chapter lemma lemma-stratify-flat-fp-lqf-universally-bounded, cited in the proof. Supplies the boundedness of the degree over the stratum, which is what makes the rank r constant there.
\item Support of a module and quasi-compact opens. Modules, definition modules-definition-support and Algebra, lemma algebra-lemma-qc-open, cited in the proof. Controls the support of the section and lets the construction be restricted to quasi-compact opens.
\item The ideal A[I] of elements annihilating I. construction used at proof lines 12966-12975. Produces the quasi-coherent subsheaf F and the map F\textasciicircum{}\{r\} -> pi\_*O\_V realising the annihilator on the base.
\end{itemize}\medskip
\ifdefined\hypertarget\hypertarget{sourcepointer-1}{}\fi
\noindent\textit{Source reference.} Original reference 00F6: focused-seeds/geometry/S9-morph-go-down-annihilators/ref-tags.txt, source lines 4--4. Navigation pointer only: the counted proof cites this statement by \textbackslash{}ref as a standing input; the deduction the counted proof has to make is carried out inside the counted span, and no part of that deduction is delegated to this pointer.\par\medskip
\ifdefined\hypertarget\hypertarget{sourcepointer-2}{}\fi
\noindent\textit{Source reference.} Original lemma 07RY: sources/stacks-more-morphisms/pinned-original-source.tex, source lines 12761--12761. Navigation pointer only: the counted proof cites this statement by \textbackslash{}ref as a standing input; the deduction the counted proof has to make is carried out inside the counted span, and no part of that deduction is delegated to this pointer.\par\medskip
\ifdefined\hypertarget\hypertarget{sourcepointer-3}{}\fi
\noindent\textit{Source reference.} Original lemma 07RZ: sources/stacks-more-morphisms/pinned-original-source.tex, source lines 12845--12845. Navigation pointer only: the counted proof cites this statement by \textbackslash{}ref as a standing input; the deduction the counted proof has to make is carried out inside the counted span, and no part of that deduction is delegated to this pointer.\par\medskip
\ifdefined\hypertarget\hypertarget{sourcepointer-4}{}\fi
\noindent\textit{Source reference.} Original reference 01AT: focused-seeds/geometry/S9-morph-go-down-annihilators/ref-tags.txt, source lines 5--5. Navigation pointer only: the counted proof cites this statement by \textbackslash{}ref as a standing input; the deduction the counted proof has to make is carried out inside the counted span, and no part of that deduction is delegated to this pointer.\par\medskip
\begin{lemma}
\label{lemma-go-down-with-annihilators}
Let $f : X \to S$ be a morphism of schemes which is flat, locally of
finite presentation, and locally quasi-finite. Let
$g \in \Gamma(X, \mathcal{O}_X)$ nonzero. Then there exist
an open $V \subset X$ such that $g|_V \not = 0$, an open
$U \subset S$ fitting into a commutative diagram
$$
\xymatrix{
V \ar[r] \ar[d]_\pi & X \ar[d]^f \\
U \ar[r] & S,
}
$$
a quasi-coherent subsheaf $\mathcal{F} \subset \mathcal{O}_U$, an integer
$r > 0$, and an injective $\mathcal{O}_U$-module map
$\mathcal{F}^{\oplus r} \to \pi_*\mathcal{O}_V$
whose image contains $g|_V$.
\end{lemma}

\begin{proof}
We may assume $X$ and $S$ affine. We obtain a filtration
$\emptyset = Z_{-1} \subset Z_0 \subset Z_1 \subset Z_2 \subset \ldots
\subset Z_n = S$ as in
Lemmas \ref{lemma-stratify-flat-fp-lqf-universally-bounded} and
\ref{lemma-stratify-flat-fp-qf}.
Let $T \subset X$ be the scheme theoretic support of the finite
$\mathcal{O}_X$-module $\Im(g : \mathcal{O}_X \to \mathcal{O}_X)$.
Note that $T$ is the support of $g$ as a section of $\mathcal{O}_X$
(Modules, Definition \ref{modules-definition-support}) and
for any open $V \subset X$ we have $g|_V \not = 0$ if and only if
$V \cap T \not = \emptyset$.
Let $r$ be the smallest integer such that $f(T) \subset Z_r$
set theoretically. Let $\xi \in T$ be a generic point of an irreducible
component of $T$ such that $f(\xi) \not \in Z_{r - 1}$ (and hence
$f(\xi) \in Z_r$). We may replace $S$ by an affine neighbourhood of
$f(\xi)$ contained in $S \setminus Z_{r - 1}$. Write $S = \Spec(A)$
and let $I = (a_1, \ldots, a_m) \subset A$ be a finitely generated ideal
such that $V(I) = Z_r$ (set theoretically, see
Algebra, Lemma \ref{algebra-lemma-qc-open}).
Since the support of $g$ is contained in $f^{-1}V(I)$ by our choice of $r$
we see that there exists an integer $N$ such that
$a_j^N g = 0$ for $j = 1, \ldots, m$. Replacing $a_j$ by $a_j^r$
we may assume that $Ig = 0$. For any $A$-module $M$ write
$M[I]$ for the $I$-torsion of $M$, i.e., $M[I] = \{m \in M \mid Im = 0\}$.
Write $X = \Spec(B)$, so $g \in B[I]$. Since $A \to B$ is flat we
see that
$$
B[I] = A[I] \otimes_A B \cong A[I] \otimes_{A/I} B/IB
$$
By our choice of $Z_r$, the $A/I$-module $B/IB$ is
finite locally free of rank $r$. Hence after replacing $S$ by
a smaller affine open neighbourhood of $f(\xi)$ we may assume
that $B/IB \cong (A/IA)^{\oplus r}$ as $A/I$-modules.
Choose a map $\psi : A^{\oplus r} \to B$ which reduces modulo $I$ to the
isomorphism of the previous sentence. Then we see that
the induced map
$$
A[I]^{\oplus r} \longrightarrow B[I]
$$
is an isomorphism. The lemma follows by taking $\mathcal{F}$ the
quasi-coherent sheaf associated to the $A$-module $A[I]$ and
the map $\mathcal{F}^{\oplus r} \to \pi_*\mathcal{O}_V$ the
one corresponding to $A[I]^{\oplus r} \subset A^{\oplus r} \to B$.
\end{proof}

\end{document}
